% Generated draft: requires mathematical and source review.
\documentclass[11pt]{article}
\usepackage[margin=25mm]{geometry}
\usepackage{fontspec}
\setmainfont{Latin Modern Roman}
\usepackage{amsmath,amssymb,mathtools,unicode-math}
\setmathfont{Latin Modern Math}
\usepackage{enumitem,needspace,xurl,hyperref}
\setlength{\parindent}{0pt}
\setlength{\parskip}{5pt}
\setlength{\emergencystretch}{3em}
\newcommand{\N}{\mathbb{N}}
\newcommand{\Z}{\mathbb{Z}}
\newcommand{\Q}{\mathbb{Q}}
\newcommand{\R}{\mathbb{R}}
\newcommand{\C}{\mathbb{C}}
\newcommand{\F}{\mathbb{F}}
\newcommand{\A}{\mathbb{A}}
\newcommand{\PP}{\mathbb{P}}
\newcommand{\E}{\mathbb{E}}
\newcommand{\eps}{\varepsilon}
\newcommand{\dd}{\,\mathrm{d}}
\begin{document}
\section*{UnsolvedMath ID: 701222}
\textbf{Source problem code:} catalogue entry
\subsection{Definitions and mathematical statement}
Let \(G\) be an arbitrary group and let \(A,B\leq G\). Define \(AB=\{ab:a\in A,\ b\in B\}\); the hypothesis is \(G=AB\). Neither finiteness of \(G\) nor normality of \(A\) or \(B\) is assumed.

For a group \(H\), a finite composition series means a chain
\[
1=H_0\triangleleft H_1\triangleleft\cdots\triangleleft H_n=H
\]
with strict inclusions and nontrivial simple factors \(H_i/H_{i-1}\) for \(1\leq i\leq n\). A nontrivial group is simple if its only normal subgroups are the identity subgroup and itself. Here \(n\) is a nonnegative integer, and \(n=0\) is permitted precisely when \(H=1\).

For this draft, interpret a \emph{principal series} as a finite chief series: a finite chain
\[
1=H_0<H_1<\cdots<H_m=H
\]
in which every \(H_i\) is normal in \(H\), and for each \(1\leq i\leq m\) there is no normal subgroup \(N\trianglelefteq H\) satisfying \(H_{i-1}<N<H_i\). Equivalently, \(H_i/H_{i-1}\) is a minimal nontrivial normal subgroup of \(H/H_{i-1}\). The empty chain is permitted for the trivial group. The supplied source does not define ``principal series'' or explicitly specify that the series are finite; these conventions require confirmation against the original source.

\textbf{Questions (V. A. Vedernikov, Kourovka Notebook Problem 16.30).} Under these conventions, are the following two universal implications valid?
\begin{enumerate}
\item For every group \(G\) and subgroups \(A,B\leq G\) with \(G=AB\), if each of \(A\) and \(B\) admits a finite composition series, then \(G\) admits a finite composition series.
\item For every group \(G\) and subgroups \(A,B\leq G\) with \(G=AB\), if each of \(A\) and \(B\) admits a finite principal series, then \(G\) admits a finite principal series.
\end{enumerate}
In the second question, the principal series of \(A\), \(B\), and \(G\) are taken relative to their respective ambient groups.

The supplied registry labels the problem open. Its discussion also reports an explicitly unverified AI claim that both implications fail. That report is a source qualification, not an established answer in this statement-only draft.
\subsection{Short English statement}
If a group is the setwise product of two subgroups, must it have a composition series when both subgroups do? Ask the same question for principal series, interpreted here as finite chief series. The source does not explicitly define its series conventions. It records an open status alongside an unverified claim of negative answers to both questions.
\subsection{Sources}
\begin{itemize}
\item Local UnsolvedMath registry, supplied record with permanent ID \(701222\), exact source problem code \texttt{catalogue entry}, title \emph{Kourovka Notebook Problem 16.30}. Proposed by V. A. Vedernikov; proposed year 2006.
\item \emph{The Kourovka Notebook (Unsolved Problems in Group Theory)}, 16th issue (2006), Problem 16.30; also attributed in the supplied record to the 21st issue edition (2026), \texttt{arXiv:1401.0300v48}. The supplied collection description names E. I. Khukhro and V. D. Mazurov as editors. These bibliographic attributions have not been independently verified here.
\item The supplied discussion cites Achyuth Jayadevan, \emph{Preprint} (2026), \url{https://zenodo.org/records/22908884}, for an explicitly unverified AI claim of negative answers for both composition and principal series. No title is supplied. The supplied status review describes the evidence as an unrefereed claim and reports that the problem bears a \(\#\) marker for an unverified AI solution claim in the October 2026 21st issue.
\item The supplied status review is dated \texttt{2026-10-16} and retains canonical status \texttt{open}. This assessment is supplied dated curation, not independently verified literature or mathematical verification.
\end{itemize}
\end{document}
